\documentclass[10pt,a4paper]{amsart}
\usepackage[margin=19mm]{geometry}
\usepackage{amsmath,amssymb,mathtools}
\usepackage[scr=rsfs,cal=euler]{mathalfa}
\usepackage{xfrac}
\usepackage[unicode,hidelinks,bookmarks=false]{hyperref}
\newcommand{\ie}{i.e.\ }
\newcommand{\cf}{cf.\ }
\newcommand{\resp}{respectively\ }
\newcommand{\Subsection}{\subsection}
\providecommand{\nobreakdash}{\nobreak\hbox{-}}
\setlength{\emergencystretch}{2em}
\multlinegap0pt
\usepackage{bm}
\usepackage{bbm}
\usepackage{dsfont}
\usepackage{xspace}
\usepackage{cancel}
\usepackage{mathtools}
\usepackage{amsthm}
\makeatletter
\@ifundefined{thm}{\newtheorem{thm}{Thm}}{}
\@ifundefined{lem}{\newtheorem{lem}[thm]{Lemma}}{}
\makeatother
\title[Optimal-Transport Stability of Inverse Point-Source Problems]{Optimal-Transport Stability of Inverse Point-Source Problems for Elliptic and Parabolic Equations}
\author{Lingyun Qiu, Shenwen Yu}
\date{Source: arXiv:2512.21821v1 (CC BY 4.0)}
\begin{document}
\maketitle

\noindent\emph{Source and licence.} Optimal-Transport Stability of Inverse Point-Source Problems for Elliptic and Parabolic Equations. Version arXiv:2512.21821v1 (CC BY 4.0), distributed under the Creative Commons Attribution 4.0 International licence, as recorded in the arXiv licence metadata for this exact version and shown on the abstract page https://arxiv.org/abs/2512.21821v1. Full rights notice: licenses/arxiv-2512.21821-CC-BY-4.0.txt.\par\smallskip

\noindent\textit{Editorial scope.} Counted unit: the Lem whose source label is \texttt{lemma:vectorb} in the article (source0.tex lines 1403--1405, source label \texttt{lemma:vectorb}), delivered together with the article's own complete original proof of it (source0.tex lines 1408--1520, one \texttt{proof} environment of 113 lines). No mathematics is written, repaired or re-derived: statement and proof are the article's own text line for line. The proof is detached from the statement in the source: the article states the result at lines 1403--1405 and prints its proof at lines 1408--1520, and the proof environment's own header names the statement, so the association is the article's own. Required context retained uncounted: CGO equations (lines 1313--1315); lemma: test function construction (lines 1328--1338); lowerbound (lines 1341--1358). Every definition, display and label of those lines is retained unchanged. Omitted from this package and not redistributed: 1--1312, 1316--1327, 1339--1340, 1359--1402, 1406--1407, 1521--1672. Nothing inside a retained excerpt is omitted or altered: every retained body is delivered exactly as the article's lines are written, so no omission receipt of any kind accompanies this package. Citations: the retained excerpts carry no citation command at all, so no bibliography entry of the article is transcribed and no reference is redistributed. Reference adaptation: none was needed and none was made; every reference inside the delivered unit resolves inside it, so no unresolved reference or citation is produced. Printed slips kept literal: no printed word, symbol, hypothesis or number of the counted statement or of its proof was altered. Layout and class: the article's own class is \texttt{amsart} with packages including newtxtext, newtxmath, amsmath, amssymb, amsfonts, amsthm, graphicx, float, booktabs, multirow, subcaption, tikz, pict2e, bm; the wrapper is the lane's pinned \texttt{amsart} editorial wrapper at 10pt on A4, which restates the theorem-like environments the retained text needs and adds the article's own macro definitions that the retained text needs (none), with extra wrapper packages (bm, bbm, dsfont, xspace, cancel, mathtools). The delivered numbering is the unit's own. Assets: no retained span contains a figure environment, a graphics-inclusion command, a PDF-page inclusion, a table or a TikZ picture; no image file is copied into this package, the package declares no asset dependency and no image file is redistributed with it. Licence: version arXiv:2512.21821v1 of this article is distributed under the Creative Commons Attribution 4.0 International licence, as recorded in the arXiv licence metadata for this exact version and shown on the abstract page https://arxiv.org/abs/2512.21821v1. A full rights notice is delivered at licenses/arxiv-2512.21821-CC-BY-4.0.txt. Characters: every retained excerpt is pure ASCII, so the delivered engine, XeLaTeX with its default Latin Modern text font, reports no missing character.
\begin{equation}\label{CGO equations}
    (\Delta + \tilde{q})\, w(x) = 0 \quad \text{in } \Omega,
\end{equation}

\begin{lem}\label{lemma: test function construction}
For every complex vector $\rho$ satisfying $\rho \cdot \rho=0$, suppose:
\begin{equation}\label{vector condition}
    |\rho|\geq C_1(\Omega,p)(1+||\tilde{q}||_{H^p(\Omega)}),
\end{equation}

where $C_1(\Omega,p)$ is a constant only related to $\Omega$ and $p$. Then we can construct a test function $v\in \mathcal{V}$ such that $v(x)=\frac{w(x)}{\sqrt{\kappa(x)}}$, where $w(x)=\mathrm{e}^{\rho \cdot x}\Big(1+\psi_\rho(x)\Big)$ satisfies \eqref{CGO equations}. Besides, the correction term $\psi_{\rho}$ belongs to $C^{2}(\Omega)$ and its norm satisfies:
\begin{equation}
    |\rho|\left\|\psi_\rho\right\|_{C^{1}(\Omega)}+\left\|\psi_\rho\right\|_{C^{2}(\Omega)} \leqslant C_1(\Omega,p) ||\tilde{q}||_{H^p(\Omega)}.
\end{equation}
\end{lem}

\begin{lem}\label{lowerbound}
    Consider a general matrix $A = (a_{l,j})_{1 \leq l, j \leq m}$. Define
    \[
    \beta(A) := \min_{\|y\|_2 = 1} \|A y\|_2, \quad y \in \mathbb{C}^m.
    \]
    Then the following estimate holds:
    \begin{equation*}
       \beta(A) \geq \min_{1 \leq j \leq m} |a_{j,j}| - \Bigg( \sum_{1 \leq j \neq l \leq m} |a_{l,j}|^2 \Bigg)^{1/2}.
    \end{equation*}
    Moreover, for any matrix $B \in \mathbb{C}^{m \times m}$, we have
    \begin{equation*}
        \beta(A + B) \geq \beta(A) - \|B\|_2, \qquad \beta(AB) \geq \beta(A)\, \beta(B).
    \end{equation*}
    In particular, if $\beta(A) > 0$, then $A$ is invertible and satisfies
    \begin{equation*}
        \|A^{-1}\|_2 \leq \frac{1}{\beta(A)}.
    \end{equation*}
\end{lem}

\begin{lem}\label{lemma:vectorb}
   For any real vector $a\in \mathbb{R}^d$, there exists another real vector $b$, such that the complex vector $\rho =a+b\mathrm{i} \in \mathbb{C}^d$ satisfying $\rho \cdot \rho=0$.
\end{lem}

\begin{proof}
For each $s_j$, we choose $\rho_j = \tilde{r} s_j + \mathrm{i} b_j$, where $b_j$ is selected according to Lemma~\ref{lemma:vectorb}. We define
\begin{equation*}
    \tilde{v}_j(x) = \frac{w_j(x)}{\sqrt{\kappa(x)\kappa(s_j)}}, \quad \text{where } w_j(x) = e^{\rho_j \cdot x}\left(1 + \psi_j(x)\right),
\end{equation*}
and $w_j$ is the CGO solution satisfying equation~\eqref{CGO equations}.

Define the diagonal matrix
\[
P = \operatorname{diag}\left( \exp\left(\frac{\tilde{r} s_j \cdot s_j}{2} \right) \frac{1}{\sqrt{\kappa(s_j)}} \right)_{1 \leq j \leq n}.
\]
Then we can write
\[
A = P (B + \delta B) P,
\]
where
\begin{align*}
    B &= \left[ \exp(\mathrm{i} b_l \cdot s_j) \exp\left(-\frac{\tilde{r} |s_j - s_l|^2}{2}\right) \right]_{1 \leq l,j \leq n}, \\
    \delta B &= \left[ \exp(\mathrm{i} b_l \cdot s_j) \exp\left(-\frac{\tilde{r} |s_j - s_l|^2}{2}\right) \psi_l(s_j) \right]_{1 \leq l,j \leq n}.
\end{align*}

We now apply Lemma~\ref{lowerbound} to estimate $\beta(B)$:
\begin{align*}
    \beta(B) 
    &\geq \min_{1 \leq j \leq n} |\exp(\mathrm{i} b_j \cdot s_j)| 
        - \left( \sum_{1 \leq j \neq l \leq n} 
        \left| \exp\left(\mathrm{i} b_l \cdot s_j\right) 
        \exp\left(-\frac{\tilde{r} |s_j - s_l|^2}{2} \right) \right|^2 \right)^{1/2} \\
    &= 1 - \left( \sum_{1 \leq j \neq l \leq n} 
        \exp\left(-\tilde{r} |s_j - s_l|^2 \right) \right)^{1/2} \\
    &\geq 1 - \sqrt{n(n-1)} \exp\left(-\frac{\tilde{r} \eta_1^2}{2} \right),
\end{align*}
where we used the uniform bound $|s_j - s_l| \geq \eta_1$ for $j \neq l$ in the last inequality. Therefore,
\begin{equation}\label{eq:betaB}
    \beta(B) \geq 1 - n \exp\left(-\frac{\tilde{r} \eta_1^2}{2} \right).
\end{equation}

Next, we estimate $\|\delta B\|_2$. Using the definition of $\delta B$ and bounding the exponential terms by $1$, we have
\begin{align*}
    \|\delta B\|_2 
    &\leq \left( \sum_{l,j=1}^n 
        \left| \exp(\mathrm{i} b_l \cdot s_j) 
        \exp\left(-\frac{\tilde{r} |s_j - s_l|^2}{2} \right) \psi_l(s_j) \right|^2 \right)^{1/2} \\
    &\leq \left( \sum_{l,j=1}^n 
        |\psi_l(s_j)|^2 \right)^{1/2}.
\end{align*}

By Lemma~\ref{lemma: test function construction}, we have the estimate
\[
|\rho_l| = \sqrt{2} \tilde{r} |s_l| \geq \sqrt{2} \tilde{r} \eta_2 
    \geq C_1\left(1 + \|\tilde{q}\|_{H^p(\Omega)} \right),
\]
which ensures that
\[
\|\psi_l\|_{C^0(\Omega)} \leq \frac{C_1 \|\tilde{q}\|_{H^p(\Omega)}}{\sqrt{2} \tilde{r} \eta_2}.
\]
Thus, we conclude that
\begin{equation}\label{eq:deltaB}
    \|\delta B\|_2 \leq n \cdot \frac{C_1 \|\tilde{q}\|_{H^p(\Omega)}}{\sqrt{2} \tilde{r} \eta_2}.
\end{equation}

Combining \eqref{eq:betaB} and \eqref{eq:deltaB}, we apply Lemma~\ref{lowerbound} to obtain
\begin{equation*}
    \beta(B + \delta B) \geq \beta(B) - \|\delta B\|_2 
    \geq 1 - n\left( \exp\left(-\frac{\tilde{r} \eta_1^2}{2} \right) + \frac{C_1 \|\tilde{q}\|_{H^p(\Omega)}}{\sqrt{2} \tilde{r} \eta_2} \right).
\end{equation*}

Since $P$ is diagonal, we have
\[
\beta(P) = \min_{1 \leq j \leq n} 
\left( \exp\left(\frac{\tilde{r} s_j \cdot s_j}{2} \right) 
\frac{1}{\sqrt{\kappa(s_j)}} \right)
\geq \frac{\exp\left(\frac{\tilde{r} \eta_2^2}{2} \right)}{\sqrt{ \|\kappa\|_{C^0(\Omega)} }}.
\]
Therefore, combining the estimates, we obtain
\begin{align*}
    \beta(A) 
    &\geq \beta(P)\, \beta(B + \delta B)\, \beta(P) \\
    &\geq \frac{\exp(\tilde{r} \eta_2^2)}{\|\kappa\|_{C^0(\Omega)}} 
    \left[ 1 - n \left( \exp\left(-\frac{\tilde{r} \eta_1^2}{2} \right) + \frac{C_1 \|\tilde{q}\|_{H^p(\Omega)}}{\sqrt{2} \tilde{r} \eta_2} \right) \right].
\end{align*}

By the definition of $\tilde{r}$, we have
\[
1 > n \left( \exp\left(-\frac{\tilde{r} \eta_1^2}{2} \right) + \frac{C_1 \|\tilde{q}\|_{H^p(\Omega)}}{\sqrt{2} \tilde{r} \eta_2} \right),
\]
which implies $\beta(A) > 0$, and hence $A$ is invertible by Lemma~\ref{lowerbound}.

It remains to estimate the $H^1$-norm of the test functions $\tilde{v}_j$. Since $\rho_j \cdot \rho_j = 0$, we observe that
\[
\tilde{r} |s_j| = |b_j| = \frac{|\rho_j|}{\sqrt{2}}.
\]
We estimate
\begin{align*}
    \|\tilde{v}_j\|_{C^1(\Omega)} 
    &\leq 3 \left\| \frac{1}{\sqrt{\kappa(x)}} \right\|_{C^1(\Omega)} 
    \left\| e^{\rho_j \cdot x} \right\|_{C^1(\Omega)} 
    \left( 1 + \|\psi_j\|_{C^1(\Omega)} \right) \\
    &\leq 3 \sqrt{2} \left\| \frac{1}{\sqrt{\kappa(x)}} \right\|_{C^1(\Omega)} 
    \tilde{r} |s_j|\, e^{\tilde{r} |s_j| R_0} 
    \left( 1 + \|\psi_j\|_{C^1(\Omega)} \right) \\
    &\leq 3 \sqrt{2} \tilde{r} |s_j| e^{\tilde{r} |s_j| R_0}\left\| \frac{1}{\sqrt{\kappa(x)}} \right\|_{C^1(\Omega)}
    \left( 1 + \frac{C_1 \|\tilde{q}\|_{H^p(\Omega)}}{\sqrt{2} \tilde{r} |s_j|} \right) \\
    &\leq 3 \sqrt{2} \tilde{r} R_0 e^{\tilde{r} R_0^2}\left\| \frac{1}{\sqrt{\kappa(x)}} \right\|_{C^1(\Omega)}
    \left( 1 + \frac{C_1 \|\tilde{q}\|_{H^p(\Omega)}}{\sqrt{2} \tilde{r} \eta_2} \right).
\end{align*}

Finally, since $\Omega$ is bounded, the $C^1$-norm controls the $H^1$-norm:
\[
\|\tilde{v}_j\|_{H^1(\Omega)} \leq \sqrt{(d + 1) |\Omega|} \cdot \|\tilde{v}_j\|_{C^1(\Omega)}.
\]
This completes the estimate for $\max_{1 \leq j \leq n} \|\tilde{v}_j\|_{H^1(\Omega)}$ and thus concludes the proof.
\end{proof}

\end{document}
